\section{The VRASection Algorithm}
\label{sec:algorithm}

\subsection{Setup and Notation}

Let $G = (V, E, w)$ be the state's census-tract adjacency graph.
Vertices $V$ are census tracts; edges $E$ connect tracts sharing a common
boundary; edge weights $w(e) > 0$ are shared boundary lengths in metres.
Vertex weights $\mathbf{p} \in \mathbb{Z}_{>0}^{|V|}$ are census tract populations.
Let $k$ be the number of congressional districts allocated to this state.
Let $\mathbf{m} \in \mathbb{R}_{\geq 0}^{|V|}$ be per-tract minority VAP
counts, where $m_v$ is the minority voting-age population of tract $v$.

\begin{definition}[Feasible bisection]
A \emph{feasible bisection} of $G$ with ratio $i:(k-i)$, $1 \leq i \leq \lfloor k/2 \rfloor$,
is a partition $(L, R)$ with $L \cup R = V$, $L \cap R = \emptyset$,
and $\sum_{v \in L} p_v / \sum_{v \in V} p_v \in [i/k \cdot (1-\delta),\; i/k \cdot (1+\delta)]$
for balance tolerance $\delta$ (congressional default: $\delta = 0.005$).
The \emph{edge-cut} of the bisection is $\mathrm{EC}(L) = \sum_{e \in \partial(L)} w(e)$
where $\partial(L)$ is the set of edges crossing the $L$/$R$ boundary.
\end{definition}

\begin{definition}[Isoperimetrically normalised edge-cut]
For ratio $i:(k-i)$, the \emph{normalised edge-cut} is:
\[
  \mathrm{EC}_\text{norm}(i) = \frac{\min_{(L,R) \text{ feasible at } i:(k-i)} \mathrm{EC}(L)}{\sqrt{\min(i,\, k-i)}}
\]
The $\sqrt{\min(i,k-i)}$ denominator is GeoSection's isoperimetric correction
\citep{b8paper}, which ensures that all ratios compete on equal footing.
\end{definition}

\subsection{Geographic Alignment Score}

\begin{definition}[Minority VAP fraction]
For a subset $S \subseteq V$, the \emph{minority VAP fraction} is:
\[
  \mathrm{MVAP}_\text{frac}(S) = \frac{\sum_{v \in S} m_v}{\sum_{v \in V} m_v}
\]
where the denominator is over the full vertex set $V$ of the current subgraph
(not just $S$). This measures what fraction of the subgraph's total minority
population lands on side $S$ of the bisection.
\end{definition}

\begin{definition}[Alignment score]
\label{def:alignment}
For a feasible bisection $(L, R)$ of $G$, the \emph{alignment score} is:
\[
  A(L, R) = \left|\mathrm{MVAP}_\text{frac}(L) - 0.5\right| \times 2
           = \left|\frac{\sum_{v \in L} m_v}{\sum_{v \in V} m_v} - 0.5\right| \times 2
\]
Equivalently, $A(L,R) = |\mathrm{MVAP}_\text{frac}(L) - \mathrm{MVAP}_\text{frac}(R)|$.
\end{definition}

The alignment score has a natural interpretation.
$A = 0$ means minority population is split exactly 50/50 between $L$ and $R$
--- the bisection is \emph{dispersive}: it divides the minority community
across both sides.
$A = 1$ means all minority population is on one side --- the bisection is
\emph{concentrating}: it places the entire minority community on one side.
A high alignment score indicates that the state has a natural geographic
boundary separating high-minority from low-minority areas, and that the
bisection at ratio $i:(k-i)$ falls along this boundary.

\subsection{Modified Selection Objective}

VRASection modifies GeoSection's ratio selection criterion by combining
the normalised edge-cut with the alignment score.

\begin{definition}[VRASection selection score]
\label{def:selection-score}
For ratio $i:(k-i)$ with best-seed bisection $(L^*, R^*)$ at that ratio, the
\emph{VRASection selection score} is:
\[
  \mathrm{score}(i) = \mathrm{EC}_\text{norm}(i)
  - w_\text{vra} \cdot A(L^*, R^*)
    \cdot \max\!\left(\mathrm{EC}_\text{norm}(i),\, 1\right)
\]
where $w_\text{vra} \in [0, 1]$ is the alignment weight (default $0.40$).
The ratio $i^*$ selected by VRASection minimises this score:
\[
  i^* = \argmin_{1 \leq i \leq \lfloor k/2 \rfloor} \mathrm{score}(i)
\]
\end{definition}

The $\max(\mathrm{EC}_\text{norm}(i), 1)$ factor ensures the alignment
bonus scales with the compactness cost: a ratio with high EC gets a
proportionally larger bonus for good alignment, so alignment can make a
moderately expensive ratio competitive with a very cheap but dispersive ratio.
For typical states where $\mathrm{EC}_\text{norm} \gg 1$ (edge cuts are in
the hundreds of kilometres), this reduces to:
\[
  \mathrm{score}(i) \approx \mathrm{EC}_\text{norm}(i) \cdot (1 - w_\text{vra} \cdot A(L^*, R^*))
\]
which has a clean interpretation: the effective compactness cost is discounted
by the alignment bonus.

\begin{proposition}[VRASection reduces to GeoSection at $w_\text{vra} = 0$]
\label{prop:geosection}
When $w_\text{vra} = 0$, $\mathrm{score}(i) = \mathrm{EC}_\text{norm}(i)$
for all $i$, and VRASection selects the same ratio as GeoSection.
\end{proposition}

\begin{proof}
Immediate from Definition~\ref{def:selection-score}: the alignment term
vanishes when $w_\text{vra} = 0$.
\end{proof}

\begin{proposition}[Pure geographic concentration at $w_\text{vra} = 1$]
\label{prop:pure-concentration}
When $w_\text{vra} = 1$, the selection score is
$\mathrm{score}(i) = \mathrm{EC}_\text{norm}(i) \cdot (1 - A(L^*, R^*))$,
and the algorithm maximises geographic alignment subject to compactness
being non-negative.
Among ratios where $A = 1$ (perfect concentration), the ratio with smallest
$\mathrm{EC}_\text{norm}$ is selected.
\end{proposition}

\begin{proof}
At $w_\text{vra} = 1$ and for $\mathrm{EC}_\text{norm} \gg 1$:
$\mathrm{score}(i) = \mathrm{EC}_\text{norm}(i) \cdot (1 - A(L^*, R^*))$.
This is minimised by maximising $A$ (since $\mathrm{EC}_\text{norm} > 0$).
When $A = 1$, score is 0 regardless of compactness, but the $\max(\cdot, 1)$
floor prevents the score from going negative: when EC is very small
(near-zero), the $\max$ ensures a positive compactness signal is retained.
\end{proof}

\begin{remark}[Proposed default $w_\text{vra} = 0.40$]
The default weight balances three considerations:
(1) \emph{Legal}: at $w_\text{vra} \leq 0.5$, compactness is the predominant
factor in the objective (see \S\ref{sec:legal-analysis});
(2) \emph{Empirical}: sensitivity analysis (B.14 pending) suggests that
$w_\text{vra} = 0.20$ often leaves alignment underutilised,
while $w_\text{vra} = 0.60$ risks over-prioritising concentration in states
with dispersed minority populations;
(3) \emph{Interpretability}: the 60/40 compactness/alignment split
is legible to non-technical audiences and can be stated plainly in
expert testimony.
\end{remark}

\subsection{Legal Analysis: Why VRASection Is Not Racial Targeting}
\label{sec:legal-analysis}

The \emph{Shaw}/\emph{Miller} racial predominance test asks:
did the legislature subordinate traditional redistricting principles
to racial considerations?
VRASection's design addresses this test on two levels.

\paragraph{What the algorithm does not do.}
VRASection does not:
\begin{itemize}
\item Set a target minority percentage for any district.
\item Require any district to achieve $> 50\,\%$ minority VAP.
\item Enumerate racial groups or treat different groups differently.
\item Use racial identity as a selection criterion at any level
  (minority VAP is a census attribute, not a racial classification).
\end{itemize}

\paragraph{What the algorithm does.}
VRASection observes the spatial distribution of a user-designated
population attribute (minority VAP) and rewards bisections that
separate high-concentration tracts from low-concentration tracts.
This is categorically a \emph{geographic} operation: the algorithm
finds a natural geographic boundary in the tract adjacency graph.
If tracts with high minority VAP were distributed uniformly across
the state, the alignment score would be near zero for all ratios,
and VRASection would select the same ratio as GeoSection.
The alignment score is non-trivial only when there is a geographic
boundary separating minority-concentrated areas from dispersed areas
--- i.e., only when Gingles Prong 1's geographic compactness requirement
is actually satisfied in the data.

\paragraph{The compactness dominance condition.}
At $w_\text{vra} < 0.5$, the compactness component of the score
($\mathrm{EC}_\text{norm}$) exceeds the maximum possible alignment discount
($w_\text{vra} \cdot 1 \cdot \mathrm{EC}_\text{norm} = w_\text{vra} \cdot \mathrm{EC}_\text{norm}$)
by a factor of at least $(1 - w_\text{vra}) / w_\text{vra} > 1$.
That is:
\[
  \text{compactness contribution} \;=\; (1 - w_\text{vra}) \cdot \mathrm{EC}_\text{norm}(i)
  \;>\; w_\text{vra} \cdot \mathrm{EC}_\text{norm}(i) \;=\; \text{alignment contribution}
\]
Geographic compactness is quantifiably the predominant factor for any
$w_\text{vra} \in (0, 0.5)$.
The proposed default $w_\text{vra} = 0.40$ gives compactness 60\,\% weight
and alignment 40\,\% weight, satisfying the \emph{Shaw} threshold with
a 20\,pp margin.

\paragraph{Contrast with MetisVra.}
MetisVra adds minority concentration as a co-equal METIS constraint.
In the ncon=2 formulation, minority VAP fraction and population balance
are simultaneous hard constraints with equal solver priority.
This likely fails \emph{Miller}'s predominance test in states where race
must override compactness to achieve the minority target --- precisely
the states where MetisVra fails empirically (AL, LA, SC).
VRASection uses race as a soft tie-breaker, bounded to 40\,\% of the
selection score, satisfying a proportionality standard rather than maximisation.

\subsection{Algorithm Pseudocode}

\begin{algorithm}[H]
\caption{VRASection($G$, $k$, $\mathbf{m}$, $N$, $\delta$, $w_\text{vra}$)}
\label{alg:vrasection}
\begin{algorithmic}[1]
\Require Graph $G = (V, E, w)$, districts $k \geq 1$, minority VAP $\mathbf{m} \in \mathbb{R}^{|V|}$,
  seeds per ratio $N$, balance tolerance $\delta$, alignment weight $w_\text{vra}$
\Ensure Partition $\{D_1, \ldots, D_k\}$ of $V$

\If{$k = 1$} \Return $\{V\}$ \EndIf

\State $M_\text{total} \gets \sum_{v \in V} m_v$
  \hfill\Comment{total minority VAP in current subgraph}

\State\Comment{\textbf{Phase 1: Ratio scan (first bisection level only)}}
\State $\mathrm{best\_score} \gets +\infty$;
   $i^* \gets 1$; $L^* \gets \emptyset$

\For{$i = 1$ to $\lfloor k/2 \rfloor$}
  \State $\mathrm{EC}_{\min} \gets +\infty$; $L_{\min} \gets \emptyset$
  \For{$s = 1$ to $N$}
    \State $(L, R) \gets \mathrm{METIS}(G,\, i:(k-i),\, \mathrm{ufactor} = 1 + \delta/k,\, \mathrm{seed}=s)$
    \If{$\mathrm{EC}(L) < \mathrm{EC}_{\min}$}
      $\mathrm{EC}_{\min} \gets \mathrm{EC}(L)$; $L_{\min} \gets L$
    \EndIf
  \EndFor
  \State $\mathrm{EC}_\text{norm}(i) \gets \mathrm{EC}_{\min} / \sqrt{\min(i, k-i)}$
    \hfill\Comment{isoperimetric normalisation}
  \If{$M_\text{total} > 0$ \textbf{and} $w_\text{vra} > 0$}
    \State $M_L \gets \sum_{v \in L_{\min}} m_v$
    \State $A_i \gets |M_L / M_\text{total} - 0.5| \times 2$
      \hfill\Comment{alignment score $\in [0, 1]$}
    \State $\mathrm{score}(i) \gets \mathrm{EC}_\text{norm}(i)
           - w_\text{vra} \cdot A_i \cdot \max(\mathrm{EC}_\text{norm}(i),\, 1)$
  \Else
    \State $\mathrm{score}(i) \gets \mathrm{EC}_\text{norm}(i)$
      \hfill\Comment{reduce to GeoSection}
  \EndIf
  \If{$\mathrm{score}(i) < \mathrm{best\_score}$}
    $\mathrm{best\_score} \gets \mathrm{score}(i)$; $i^* \gets i$; $L^* \gets L_{\min}$
  \EndIf
\EndFor

\State\Comment{\textbf{Phase 2: Recursion with standard GeoSection (no VRA alignment)}}
\State $V_L \gets L^*$;\ $V_R \gets V \setminus L^*$
\State $G_L \gets G[V_L]$;\ $G_R \gets G[V_R]$
  \hfill\Comment{induced subgraphs}

\State \Return $\mathrm{GeoSection}(G_L, i^*, N, \delta)$
         $\cup$ $\mathrm{GeoSection}(G_R, k - i^*, N, \delta)$
\end{algorithmic}
\end{algorithm}

\paragraph{Design rationale: VRA alignment only at level 1.}
VRASection applies the alignment bonus only at the first bisection level.
Recursive calls use standard GeoSection (\texttt{ncon=1}, $w_\text{vra}=0$).
This choice is deliberate and has both algorithmic and legal motivations.

Algorithmically, the first bisection level has the most geographic information:
the full state's minority population distribution is visible.
At deeper recursion levels, the subgraph has already been split along a
geographic boundary; the remaining minority population is either already
concentrated (GeoSection will naturally preserve it) or dispersed (no
further alignment is available).

Legally, applying alignment at every level would make race increasingly
influential across many bisection decisions, accumulating into racial
predominance across the full bisection tree.
A single first-level alignment observation, bounded to 40\,\% of the
selection score, is categorically different from race influencing every
split at every level.
The algorithm can state: ``race was observed once, at the highest geographic
scale, as a minor factor; all subsequent decisions were geographically
determined.''

\paragraph{Computational cost.}
VRASection adds zero METIS calls compared to GeoSection: the alignment
score is computed from the bisection assignments already produced by the
ratio scan, not from additional METIS runs.
The cost is one pass over the vertex set to accumulate minority VAP
sums, which is $O(|V|)$ per ratio.
Total computational overhead versus GeoSection:
$O(\lfloor k/2 \rfloor \cdot |V|)$ additional arithmetic operations
at the first bisection level, negligible compared to the METIS calls.
